% Part: lambda-calculus
% Chapter: lambda-definability
% Section: partial-recursive-functions

\documentclass[../../../include/open-logic-section]{subfiles}

\begin{document}

\olfileid{lam}{ldf}{par}
\olsection{الدوال العودية الجزئية \usetoken{S}{lambda definable}}

الدوال العودية الجزئية هي الدوال التي نحصل عليها من الدوال الأساسية
بالتركيب والعودية البدائية والتصغير غير المحدود. وهي تختلف عن الدوال
العودية العامة في أن الدوال المستعملة في البحث غير المحدود لا يشترط
أن تكون نظامية. وعدم اشتراط النظامية يعني أن الدوال المعرّفة بالتصغير
قد تكون غير معرّفة في بعض الأحيان.

قد يبدو للوهلة الأولى أن الطرق نفسها المستعملة لإثبات أن جميع الدوال
العودية العامة (الكلية) !!{lambda definable} يمكن استعمالها لإثبات أن
جميع الدوال العودية الجزئية !!{lambda definable}. فمثلًا، يكون تركيب
$f$ مع~$g$ !!{lambda defined} بالحد~$\lambd[x][F (G x)]$ إذا كانت
$f$ و$g$ !!{lambda defined} بالحدين~$F$ و$G$ على الترتيب. لكن هذا
يوقعنا في مشكلة عندما تكون الدوال جزئية. فإذا كانت~$g(x)$ غير معرّفة،
كان معنى ذلك أن~$G x$ لا صورة سوية له. ويعني هذا في معظم الحالات أن
$F (G x)$ لا صورة سوية له أيضًا، وهو ما نريده. لكن انظر إلى الحالة
التي يكون فيها~$F$ هو~$\lambd[x][\lambd[y][y]]$؛ فعندئذ تكون لـ
$F (G x)$ صورة سوية هي~$\lambd[y][y]$.

هذه المشكلة ليست مستعصية، وتوجد طرائق !!{lambda define} بها جميع
الدوال العودية الجزئية بحيث تمثل القيم غير المعرّفة بحدود لا صورة
سوية لها. غير أن هذه الطرائق أعقد قليلًا وأقل بداهة من النهج الذي
اتبعناه للدوال العودية العامة. ونسجل المبرهنة هنا من دون برهان:

\begin{thm}
  جميع الدوال العودية الجزئية !!{lambda definable}.
\end{thm}

\end{document}
